\documentclass[11pt, a4paper]{article}

\usepackage[margin=1in]{geometry}
\usepackage[T1]{fontenc}
\usepackage{lmodern}
\usepackage{microtype}
\usepackage{amsmath, amssymb, amsthm}
\usepackage{mathtools}
\usepackage{booktabs}
\usepackage{tabularx}
\usepackage{enumitem}
\usepackage{xcolor}
\usepackage{listings}
\usepackage[hidelinks]{hyperref}
\usepackage{url}
\usepackage{parskip}

% ── Listings style ────────────────────────────────────────────────
\definecolor{codebg}{HTML}{F5F5F7}
\definecolor{codekw}{HTML}{0F172A}
\definecolor{codecm}{HTML}{6B7280}

\lstdefinestyle{mono}{
  basicstyle=\small\ttfamily,
  backgroundcolor=\color{codebg},
  commentstyle=\color{codecm}\itshape,
  keywordstyle=\color{codekw}\bfseries,
  frame=single,
  framerule=0pt,
  xleftmargin=10pt,
  xrightmargin=10pt,
  aboveskip=6pt,
  belowskip=6pt,
  breaklines=true,
  showstringspaces=false,
}
\lstset{style=mono}

% ── Theorem environments ─────────────────────────────────────────
\newtheorem{theorem}{Theorem}
\newtheorem{definition}[theorem]{Definition}
\newtheorem{proposition}[theorem]{Proposition}

% ── Title ────────────────────────────────────────────────────────
\title{\textbf{Veritas:} A Lean-Backed Pre-Trade Verifier \\ for Agent-Native Finance}
\author{The Veritas Project \\ \texttt{\small github.com/lizjq66/Veritas}}
\date{April 2026}

\begin{document}
\maketitle

\begin{abstract}
\noindent
Automated trading agents increasingly act on capital through exchange APIs
at rates that outpace manual oversight. The public discussion frames this as
a \emph{prediction quality} problem and invests accordingly in larger models
and denser backtests. We argue that the more load-bearing failure mode of
current agent stacks is \emph{pre-trade approval}: agents take trades whose
operators would have intervened had they been given the chance, and do so
without producing a legible account of what they are about to do.

We present \textbf{Veritas}, a pre-trade verification system with a small
fixed contract. Given a trade proposal, account constraints, and existing
portfolio, Veritas returns a structured certificate indicating
\textsc{approve} / \textsc{resize} / \textsc{reject} for each of three
orthogonal gates: (1)~signal consistency, dispatched over a policy
registry; (2)~strategy-constraint compatibility, bounded by a
reliability-adjusted position-size ceiling; (3)~portfolio interference,
measured with correlation weighting across assets. Gate logic is written
in Lean~4~\cite{lean4} with closed soundness proofs; every theorem
depends only on Lean's three foundational axioms — no IEEE~754 rounding
assumptions, no Veritas-specific axioms of any kind — by virtue of
carrying out all decision-path arithmetic over exact rationals via
Mathlib~\cite{mathlib}. Python is transport. Veritas is not a trading
agent; it is the thing agents call before moving capital.
\end{abstract}

% ── 1. Introduction ──────────────────────────────────────────────
\section{Introduction}
\label{sec:intro}

The first visible failure of automated trading agents is that they lose
money. The first visible fix, therefore, is that they should predict better.
Every LLM-trading-agent release ships with a benchmark bar chart.

Empirically, though, the first several generations of agent stacks fail in
a recognizable second way: they take trades their operators would have
blocked if they had been given the chance. An agent decides to lever up a
correlated basket one minute before a scheduled Fed release. An agent opens
a ten-percent notional position on a pair whose liquidity it has not traded
before. An agent doubles down on a rejected thesis because the reviewer
prompt does not surface the reject.

These are not prediction failures. They are \emph{approval} failures. The
agent knew what it wanted to do; nothing in the system forced it to
articulate the claim, check the claim against an envelope, and account for
existing exposure before committing.

The industry response so far has been either (a)~add another LLM to critique
the first, or (b)~bolt post-trade risk controls on top of an opaque
decision loop. The first restates the problem in a larger system; the
second is reactive by construction --- it learns about a breach after the
money has moved.

Between intent and execution, there should be a \emph{pre-trade} gate that
is (i)~narrower than the agent, (ii)~legibly auditable, and (iii)~not
produced by the same process that produced the trade. This is the space
Veritas occupies.

\paragraph{Thesis.}
We argue that \emph{pre-trade verification}, not prediction quality, is
the load-bearing safety property of agent-native trading stacks, and
that the right instantiation is a narrow verifier with a
formally-verified kernel --- not a better trader with safety bolted on.
The remainder of this paper develops this claim and reports on a working
implementation.

\paragraph{Contributions.}
\begin{itemize}[leftmargin=1.5em, itemsep=2pt, topsep=2pt]
\item A three-gate decomposition of pre-trade verification, with each
      gate carrying its own soundness theorem (Section~\ref{sec:gates},
      Section~\ref{sec:formal}).
\item A Lean~4 implementation of the gates plus their supporting
      finance and learning layers, proved against exact rational
      arithmetic. Every theorem --- five gate-layer, two on the Kelly
      fraction, five on the position-sizing function, one on exit
      exhaustiveness, two on reliability updates, three on multi-
      assumption aggregation --- depends only on Lean's foundational
      \texttt{propext}, \texttt{Classical.choice}, and \texttt{Quot.sound}
      axioms. No Veritas-specific axiom is introduced.
\item A uniform HTTP / MCP~\cite{mcp} / Python / CLI contract over a
      single kernel, covered by 114 tests that exercise the compiled
      verifier with no mocks.
\item CI-enforced boundary invariants preventing the Python layer from
      reintroducing decision logic or minting approvals outside the
      kernel.
\end{itemize}

% ── 2. The verifier boundary ─────────────────────────────────────
\section{The Verifier Boundary}
\label{sec:boundary}

The formal-methods community has decades of experience with tools that sit
between a producer and a consumer and discharge obligations on the
producer's behalf: model checkers, static analyzers, type systems, SMT
solvers. The design pattern is consistent: small, fixed input contract;
small, structured output; no shared state with the caller; decision kernel
isolated from I/O.

A trading-agent verifier deserves the same shape. We adopt it directly:

\begin{itemize}[leftmargin=1.5em, itemsep=2pt, topsep=2pt]
\item \textbf{Small fixed contract.} A \texttt{TradeProposal}, an
      \texttt{AccountConstraints}, a \texttt{Portfolio}. Every field is a
      simple value. No callbacks, no futures, no handles.
\item \textbf{Small structured output.} A \texttt{Certificate} of three
      per-gate verdicts (\textsc{approve} / \textsc{resize} /
      \textsc{reject}) plus reason codes and attached assumptions. No
      exceptions.
\item \textbf{No shared state.} Veritas is a pure function: no session,
      no token, no subscription.
\item \textbf{Isolated kernel.} Gate logic lives in a Lean~4 kernel
      compiled to a native binary. The Python adapters that speak HTTP,
      MCP, and SQLite contain no gate decisions.
\end{itemize}

A verifier with this shape composes well with existing agents. A caller
can adopt it by adding one call per proposal, with no architectural
change. The verifier does not require access to the caller's alpha, model
weights, or prompts; it only needs the output of the decision the caller
is about to make.

% ── 3. The three gates ──────────────────────────────────────────
\section{The Three Gates}
\label{sec:gates}

We decompose the approval obligation into three orthogonal questions.
Each is a gate returning one of \textsc{approve}, \textsc{resize}, or
\textsc{reject}.

\subsection{Gate~1 --- Signal Consistency}
Given the proposed trade and the market context the caller claims to be
acting under, does the proposal cohere? Gate~1 dispatches over a
\emph{policy registry} --- a list of named strategies, each with its own
decider and assumption extractor. For every submitted proposal, Gate~1
runs every registered strategy's decider on the context, collects the
set of firing signals, and verifies four conditions:
(i)~at least one strategy fires, (ii)~all firing strategies agree on
direction (\texttt{MutuallyConsistent}), (iii)~the proposal's direction
matches them, and (iv)~the union of attached assumptions is non-empty.
On approval Gate~1 returns the union, so downstream auditors see every
bet the caller is implicitly committing to.

The non-triviality of Gate~1 depends on the registry containing at
least two strategies that can meaningfully disagree. The reference
registry ships with \texttt{funding\_reversion} and
\texttt{basis\_reversion}, which produce opposite-direction signals on
the same market context under conditions of long-side pressure
(positive funding rate plus perp-rich-vs-spot). That disagreement is
the case Gate~1's consistency check is designed to detect and reject.
Callers may extend the registry with their own strategies without
touching Gate~1's dispatcher.

\subsection{Gate~2 --- Strategy-Constraint Compatibility}
Given the caller's account envelope (equity, reliability of the attached
assumption(s), sample size, leverage cap, stop-loss), is the requested
notional inside the ceiling? Gate~2 returns the ceiling when the caller
is over it (\textsc{resize}), rather than silently approving at the larger
number. Gate~2 rejects when no non-zero size is admissible --- below the
reliability threshold, with non-positive leverage, or with non-positive
notional. This is the gate whose formal content is strongest: any value
it approves or resizes to is bounded by
\texttt{Finance.calculatePositionSize}, whose own five theorems
characterize the ceiling. When Gate~1 attaches multiple assumptions to an
approval, a learning-layer function \texttt{aggregateReliability} reduces
the attached list to a single conservative pair
\texttt{(reliability, sampleSize)} that Gate~2 then consumes; a single
under-sampled assumption is enough to push the proposal into
exploration-phase sizing (Section~\ref{sec:formal}).

\subsection{Gate~3 --- Portfolio Interference}
Given the caller's existing positions and the (possibly resized)
proposal, does the combined exposure violate portfolio-wide
constraints? Gate~3 measures exposure with a \emph{correlation
weighting}: each existing position contributes to the proposal's
risk bucket in proportion to the absolute value of its correlation
coefficient with the proposal's asset. Same-asset positions count at
$1$, unknown cross-asset pairs count at $0$, and explicit entries
override both defaults. Direction conflicts are restricted to
same-asset positions, so an opposite-direction position on an
uncorrelated asset (a potential hedge) is not falsely flagged as a
conflict.

\medskip

\noindent The gates execute in order. Gate~2 sees the original proposal;
Gate~3 sees the Gate~2 output. If any gate rejects, downstream gates
receive an \texttt{upstream\_gate\_rejected} marker and the final
notional is zero.

\begin{table}[h]
\small
\begin{tabularx}{\linewidth}{@{}lXl@{}}
\toprule
Gate & Question & Output \\
\midrule
1. Signal consistency
  & Registered strategies agree on a direction the proposal matches?
  & approve / reject \\
2. Constraint compat.
  & Size fits the reliability-adjusted ceiling?
  & approve / resize / reject \\
3. Portfolio interference
  & Correlation-weighted exposure respects the cap?
  & approve / resize / reject \\
\bottomrule
\end{tabularx}
\caption{The three gates.}
\label{tab:gates}
\end{table}

% ── 4. Lean-backed kernel ───────────────────────────────────────
\section{Lean-Backed Kernel}
\label{sec:lean}

The design leans on Lean~4 for the kernel. Two properties matter.

\paragraph{Compile-checked API surface.} Veritas encodes properties the
type system would otherwise leave as conventions. Exit reasons are a sum
type with three constructors (\texttt{AssumptionMet},
\texttt{AssumptionBroke}, \texttt{StopLoss}); \texttt{ReliabilityStats}
carries \texttt{wins} $\leq$ \texttt{total} at the type level, so the
constructor cannot be called with inconsistent counts; \texttt{Verdict}
has three tags, each carrying the data needed for that tag and nothing
more.

\paragraph{Theorem-backed numeric behavior.}
\texttt{Finance.PositionSizing} ships five theorems: non-negativity,
a 25\% equity cap, zero-at-no-edge, monotonicity in reliability, and
an exploration cap at 1\%.
\texttt{Finance.Kelly} ships non-negativity and monotonicity.
\texttt{Strategy.ExitLogic} ships \texttt{exitReason\_exhaustive}.
\texttt{Learning.Reliability} ships bounded-in-$[0,1]$ and
monotone-on-wins, plus three theorems pinning down the multi-
assumption aggregation. These eliminate categories of silent
miscount; a reviewer has a small set of facts about which to argue.

\paragraph{Arithmetic basis.} All decision-path arithmetic runs on exact
\texttt{Rat} (rationals) backed by Mathlib~\cite{mathlib}. Every
theorem in this paper depends only on Lean's three foundational axioms:
\texttt{propext}, \texttt{Classical.choice}, and \texttt{Quot.sound} ---
the same basis Mathlib itself rests on. There is no Veritas-specific
axiom, no rounding assumption, no Float conversion primitive; a reviewer
who trusts Mathlib (as any non-trivial Lean project does) inherits
Veritas's numeric claims at no additional cost of trust. Decimal input
from the CLI is parsed directly into \texttt{Rat} via Lean's JSON
number parser; \texttt{Rat} results are converted to \texttt{Float}
only at the JSON output boundary for wire compatibility.

% ── 5. Formal guarantees ────────────────────────────────────────
\section{Gate-Layer Soundness Theorems}
\label{sec:formal}

Each gate file exports a soundness theorem stating what an \textsc{approve}
(or \textsc{resize}) verdict from that gate \emph{means}, independent of
the underlying strategy and finance modules. The theorems are summarized
in Table~\ref{tab:thms}.

\begin{table}[h]
\footnotesize
\begin{tabularx}{\linewidth}{@{}lX@{}}
\toprule
Theorem & Statement (informal) \\
\midrule
\texttt{verifySignal\_approve\_}
\texttt{implies\_consistent}
  & Gate~1 approve $\Rightarrow$ the proposal is signal-consistent
    against the policy registry: at least one strategy fires, the
    firing strategies agree on direction, the proposal matches that
    direction, and the attached assumption list is non-empty. \\[4pt]
\texttt{checkConstraints\_approve\_}
\texttt{within\_ceiling}
  & Gate~2 approve $\Rightarrow$ $p.\texttt{notional} \leq \texttt{ceiling}$. \\[4pt]
\texttt{checkConstraints\_resize\_}
\texttt{respects\_ceiling}
  & Gate~2 resize~$n$ $\Rightarrow$ $n \leq \texttt{ceiling}$. \\[4pt]
\texttt{checkPortfolio\_approve\_}
\texttt{respects\_cap}
  & Gate~3 approve $\Rightarrow$ correlation-adjusted exposure plus
    the proposal's notional stays within the portfolio cap. \\[4pt]
\texttt{certificate\_soundness}
  & \texttt{approves} true $\Rightarrow$ Gate~1 consistent and Gate~2/3
    non-reject. \\
\bottomrule
\end{tabularx}
\caption{Gate-layer soundness theorems in \texttt{Veritas/Gates/}.}
\label{tab:thms}
\end{table}

When Gate~1 approves a proposal against the policy registry it attaches
the union of assumptions from every firing strategy. Gate~2, in
contrast, consumes a scalar \texttt{(reliability, sampleSize)} pair.
The learning layer bridges the two with an \texttt{aggregateReliability}
function that takes the element-wise minimum across its inputs. Four
supporting theorems pin down its shape:

\begin{itemize}[leftmargin=1.5em, itemsep=1pt, topsep=2pt]
\item \texttt{aggregateReliability\_empty} --- empty list returns the
      no-data default $(1/2, 0)$.
\item \texttt{aggregateReliability\_singleton} --- single-element
      lists are pass-through.
\item \texttt{aggregateReliability\_sampleSize\_le\_head} --- the
      aggregate sample size never exceeds any input's total, so a
      single under-sampled assumption is enough to push Gate~2 into
      exploration-phase sizing.
\item \texttt{aggregateReliability\_score\_le\_each} --- the
      aggregate reliability is $\leq$ every input's reliability; the
      property that justifies calling the aggregation "conservative".
\end{itemize}

\noindent The last of these is unprovable under IEEE~754 arithmetic;
it becomes routine under exact rationals.

As an example, Gate~2's ceiling theorem states:
\begin{theorem}[\texttt{checkConstraints\_approve\_within\_ceiling}]
\label{thm:gate2}
For every proposal $p$ and account constraints $c$, if
$\texttt{checkConstraints}(p, c) = \textsc{approve}$, then
\[
  p.\texttt{notionalUsd} \;\leq\;
  \texttt{calculatePositionSize}(c.\texttt{equity},\, c.\texttt{reliability},\, c.\texttt{sampleSize}).
\]
\end{theorem}
\noindent The right-hand side is in turn bounded by the five
\texttt{positionSize\_*} theorems: non-negative, $\leq \tfrac{1}{4} \cdot
c.\texttt{equity}$, zero when reliability $\leq \tfrac{1}{2}$
post-exploration, monotone in reliability, fixed at $\tfrac{1}{100}
\cdot c.\texttt{equity}$ during the first ten samples.

The combined soundness theorem connects the chain:
\begin{theorem}[\texttt{certificate\_soundness}]
For every proposal $p$, account constraints $c$, and portfolio
$\text{port}$, let
$\texttt{cert} \coloneqq \texttt{emitCertificate}(p, c, \text{port})$.
If $\texttt{cert}.\texttt{approves} = \text{true}$, then
\begin{enumerate}[leftmargin=2em, itemsep=1pt, topsep=2pt]
\item $\texttt{signalConsistent}(p)$ holds;
\item the certificate's \texttt{gate2} verdict is not a rejection;
\item the certificate's \texttt{gate3} verdict is not a rejection.
\end{enumerate}
\end{theorem}
\noindent Numeric bounds for Gates~2 and~3 follow by instantiating the
per-gate theorems on the certificate's actual verdicts.

\paragraph{Python cannot bypass the kernel.} The Python layer is
constrained by three CI-enforced invariants:
\begin{enumerate}[leftmargin=2em, itemsep=1pt, topsep=2pt]
\item No Python file branches on Veritas decision types
      (\texttt{Signal}, \texttt{ExitDecision}, \texttt{PositionSize}).
\item No Python file constructs \texttt{Verdict} values outside
      \texttt{python/schemas.py} (the dataclass) or
      \texttt{python/verifier.py} (deserialization from the kernel).
\item No Python file invokes \texttt{veritas-core} via
      \texttt{subprocess} outside \texttt{python/bridge.py}.
\end{enumerate}
These are checked by \texttt{tests/test\_bypass\_invariant.py} and
\texttt{tests/test\_loop.py}.

% ── 6. Claim discipline ─────────────────────────────────────────
\section{Claim Discipline}
\label{sec:claims}

We want to be explicit about the scope of what Veritas publishes.

\paragraph{What Veritas verifies.} That the approved notional returned by
the Lean kernel satisfies the theorems of Section~\ref{sec:formal}. That
every exit classified by \texttt{Strategy.ExitLogic} is one of three
exhaustive reasons. That reliability updates stay in $[0,1]$ and are
monotone on wins. That the multi-assumption aggregation is
element-wise minimum, and therefore never less conservative than any
input. That each gate in \texttt{Veritas/Gates/} executes the pure
function its specification describes. All of the above is closed under
Lean's three foundational axioms (Section~\ref{sec:lean}).

\paragraph{What Veritas does not verify.}
\begin{itemize}[leftmargin=1.5em, itemsep=1pt, topsep=2pt]
\item The correctness of the caller's alpha. An adversarial proposal
      with a fabricated funding rate is checked against policy but not
      against a ground-truth data feed.
\item The profitability of any bundled reference strategy. Veritas
      verifies that a strategy is applied consistently; it does not
      claim the strategy makes money.
\item Anything about a specific execution venue. After the certificate
      is issued, what the caller does is the caller's problem.
\item Correctness of the Python transport. An adapter bug could fail
      to deliver a certificate; no theorem prevents that. Callers are
      expected to treat missing certificates as rejections.
\end{itemize}

% ── 7. Related work ─────────────────────────────────────────────
\section{Related Work}
\label{sec:related}

Institutional trading desks routinely run pre-trade risk checks; the
novelty here is not the category, it is the combination of (i)~an
agent-native contract, (ii)~a Lean-backed kernel grounded in exact
rational arithmetic, and (iii)~public, legible certificates with
machine-readable reason codes.

Formal-methods work in finance has historically concentrated on
smart-contract verification. Tools such as Certora~\cite{certora} and
the Move Prover~\cite{moveprover} formalize on-chain logic. Veritas
formalizes \emph{off-chain policy} --- the step that decides whether an
order should be sent at all --- and leaves the on-chain side to existing
tools. Mathlib~\cite{mathlib} supplies the formal mathematics library
on which Veritas's rational arithmetic layer rests.

Agent-evaluation work --- benchmarks such as
AgentBench~\cite{agentbench}, the Inspect framework for red-teaming
agents, and domain-specific suites for trading agents --- focuses on
pre-deployment screening. Veritas is orthogonal: it is a run-time gate,
not a pre-deployment evaluation.

The Model Context Protocol~\cite{mcp} provides a uniform way for agents
to call tools; Veritas ships an MCP tool (\texttt{verify\_proposal})
alongside the HTTP and Python surfaces.

% ── 8. Implementation ───────────────────────────────────────────
\section{Implementation}
\label{sec:impl}

\paragraph{Deliverables.} The \texttt{Verifier} Python surface, the
\texttt{POST /verify/*} HTTP endpoints, the MCP
\texttt{verify\_proposal} tool, and the \texttt{veritas-core} CLI all
produce the same certificate for the same input, enforced by 114
passing tests with no mocks on the kernel. \texttt{lake build}
produces the compiled kernel binary with zero \texttt{sorry}. The
decision-path arithmetic runs on exact \texttt{Rat}; there are no
Veritas-specific axioms, and \texttt{\#print axioms} on every
gate-level and underlying theorem reports only \texttt{propext},
\texttt{Classical.choice}, and \texttt{Quot.sound}.

\paragraph{Example runner.} The repository ships a reference
trading agent (funding-reversion policy, Hyperliquid adapters) that
calls Veritas on every tick. It is a demonstration of a caller, not
the product. Its adapters (\texttt{HyperliquidObserver},
\texttt{HyperliquidExecutor}) are reference integrations; the
verifier's correctness does not depend on their network behaviour.

\paragraph{Interactive playground.} A browser-visible playground at
the repository's \texttt{/} route renders the three-gate verdicts
visually and fires the \texttt{POST /verify/proposal} round-trip on a
battery of preset scenarios (clean approve, direction conflict,
strategies-contradict reject, oversize $\rightarrow$ resize, no-edge
reject, portfolio conflict). It exists for human evaluators; agents
bypass it and use the HTTP contract directly.

% ── 9. Conclusion ───────────────────────────────────────────────
\section{Conclusion}
\label{sec:conclusion}

The question a trading agent should answer before touching capital is
not \emph{``what do you want to do?''} but \emph{``what do you want to
do, under what assumption, inside what envelope, against what you
already hold, and why?''} Veritas exists because nothing else in the
agent-native stack asks that question in a way the caller can be
forced through and the reviewer can audit. We would rather have a
small, legible gate than a large, opaque trader. The instantiation
described here is the smallest such gate we could ship while keeping
all three questions on the critical path; its formal and engineering
scope is the evidence that the thesis of Section~\ref{sec:intro} is
buildable, not merely arguable.

% ── References ───────────────────────────────────────────────────
\begin{thebibliography}{9}

\bibitem{lean4}
Leonardo de Moura and Sebastian Ullrich.
\emph{The Lean 4 Theorem Prover and Programming Language}.
Proceedings of CADE 28, 2021.

\bibitem{mathlib}
The Mathlib Community.
\emph{The Lean Mathematical Library}.
Proceedings of CPP 2020, ACM, 2020.

\bibitem{mcp}
Anthropic.
\emph{Model Context Protocol Specification}, 2024.
\url{https://modelcontextprotocol.io}

\bibitem{certora}
Certora.
\emph{The Certora Prover: formal verification for smart contracts}.
\url{https://www.certora.com}

\bibitem{moveprover}
David Dill, Wolfgang Grieskamp, Junkil Park, Shaz Qadeer, Meng Xu, Emma Zhong.
\emph{Fast and Reliable Formal Verification of Smart Contracts with the Move Prover}.
TACAS 2022.

\bibitem{agentbench}
Xiao Liu, Hao Yu, Hanchen Zhang, Yifan Xu, Xuanyu Lei, Hanyu Lai, Yu Gu,
Hangliang Ding, Kaiwen Men, Kejuan Yang, Shudan Zhang, Xiang Deng, Aohan Zeng,
Zhengxiao Du, Chenhui Zhang, Sheng Shen, Tianjun Zhang, Yu Su, Huan Sun,
Minlie Huang, Yuxiao Dong, Jie Tang.
\emph{AgentBench: Evaluating LLMs as Agents}.
ICLR 2024.

\end{thebibliography}

\end{document}
